\newpage
\section{Wstęp i cel pracy}
Współczesny rozwój oprogramowania staje się coraz bardziej skomplikowany, a presja na szybkie dostarczanie wysokiej jakości produktów wymaga od programistów zarówno dużego nakładu pracy, jak i stosowania narzędzi zwiększających efektywność. W ostatnich latach firmy takie jak OpenAI, Anthropic czy Meta udostępniły duże modele językowe (Large Language Models, LLM), a wraz z nimi powstały narzędzia i biblioteki (zarówno komercyjne, jak i otwarte), umożliwiające budowę systemów wspierających proces wytwarzania oprogramowania.

Obecne badania nad modelami językowymi w kontekście rozwoju oprogramowania, takie jak SWE-bench, przedstawiają test porównawczy dla generowania kodu na podstawie języka naturalnego dla różnych modeli i narzędzi. Wykazano, że modele językowe potrafią generować kod imitujący rzeczywiste poprawki błędów w rzeczywistych repozytoriach, jednak nawet najnowsze modele wykazały tylko 67,6\% skuteczności dla wszystkich badanych problemów \cite{jimenez2024swebenchlanguagemodelsresolve}. Istnieją również bardziej wymagające rankingi ocen jak LiveCodeBench, które badają nie tylko samą generację, lecz również umiejętności naprawy i testowania kodu \cite{jain2024livecodebench}. W obszarze produktywności dużą uwagę przyciągnęły badania nad GitHub Copilot, które w kontrolowanych eksperymentach skróciły czas realizacji zadań o ponad 50\% \cite{peng2023impact}.

W kontekście rozwoju oprogramowania pojawiły się systemy agentowe, takie jak SWE-agent, które dzięki dedykowanemu interfejsowi (Agent-Computer Interface) znacząco poprawiły wyniki na SWE-bench \cite{yang2024sweagent}. Inny nurt rozwoju stanowi MetaGPT, czyli podejście wieloagentowe z rolami i procedurami (SOP), które ogranicza błędy kompozycji i weryfikuje pośrednie wyniki \cite{hong2023metagpt}.

Mimo tych osiągnięć, badania nadal wskazują na brak kompleksowych analiz wpływu agentów AI na cały cykl rozwoju oprogramowania w projektach biznesowych, w tym na jakość, koszt budowy i utrzymywania \cite{cao2025buildbenchmarkrevisiting274}. Dotychczasowe prace koncentrują się głównie na benchmarkach akademickich, eksperymentach w środowisku laboratoryjnym lub wąskich analizach produktywności. Brakuje badań empirycznych w kontekście rzeczywistych projektów biznesowych, gdzie poza czasem realizacji istotne są również takie aspekty jak wskaźnik wad, czytelność kodu czy koszty utrzymania.

\subsection{Cel pracy}
Celem niniejszej pracy jest empiryczna ocena wpływu agentów AI opartych na LLM na proces tworzenia oprogramowania w realistycznym projekcie biznesowym. W szczególności porównanie różnych trybów pracy: tradycyjnego, wspieranego przez pojedyncze LLM oraz opartego o systemy wieloagentowe.

W pomiarach uwzględniono wszechstronne metryki - nie tylko produktywności (czas wykonania), ale także jakości kodu i łatwości w utrzymaniu (złożoność, czytelność) oraz kosztu obliczeniowego (liczba tokenów, czas obliczeń). Weryfikacja wpływu architektury interakcji - sprawdzenie, czy wprowadzenie elementów inspirowanych interfejsem akcji (ACI) oraz procedurami SOP wpływa na efektywność agentów, zgodnie z wynikami raportowanymi w literaturze \cite{yang2024sweagent}.

\subsection{Zakres i struktura pracy}
Praca składa się z kilku zasadniczych części. Pierwsza z nich to przegląd literatury, omówienie aspektów działania modeli językowych, agentów AI oraz ewolucji benchmarków dla LLM (HumanEval, LiveCodeBench), oceny narzędzi wspierających programistów (Copilot), architektur agentowych (SWE-agent, MetaGPT) oraz aktualnych luk badawczych. Następnie przedstawiono użyte narzędzia. Opisano zaprojektowane aplikacje biznesowe oraz dwa warianty podejścia: z pojedynczym agentem i systemem wieloagentowym z elementami ACI/SOP. Metodologia badań – definicja zadań (implementacja, poprawa defektów), metryk (czas, jakość, utrzymywalność, koszt). Zaprezentowano wyniki oraz porównano je do wyników znanych z badań akademickich. Na koniec dokonano interpretacji wyników w kontekście praktycznym oraz omówiono ograniczenia.